\begin{figure}[ht]
\centering
\begin{tikzpicture}[x=1.15cm, y=1.15cm, font=\small]

% Axes.
\draw[->, thick] (-0.3, 0) -- (5.6, 0)
  node[right, font=\scriptsize] {$x_{1}$};
\draw[->, thick] (0, -0.3) -- (0, 4.8)
  node[above, font=\scriptsize] {$x_{2}$};

% Feasible set: a convex polygon (a box-like region with a clipped corner).
\fill[engblue!8]
  (1.0, 0.6) -- (4.6, 0.6) -- (4.6, 2.6) -- (3.2, 3.8)
  -- (1.0, 3.8) -- cycle;
\draw[engblue, thick]
  (1.0, 0.6) -- (4.6, 0.6) -- (4.6, 2.6) -- (3.2, 3.8)
  -- (1.0, 3.8) -- cycle;
\node[engblue, anchor=west, font=\scriptsize] at (1.15, 1.15)
  {feasible set $\mathcal{C}$};

% Unconstrained descent target outside the set (upper right).
\fill[enggray] (5.2, 4.2) circle (1.6pt);
\node[enggray, anchor=south, font=\scriptsize] at (5.2, 4.3)
  {$\vect{x}_{k} - \eta\grad f$};

% Current iterate on the boundary.
\fill[engred] (3.4, 3.65) circle (1.8pt);
\node[engred, anchor=north east, font=\scriptsize] at (3.35, 3.6)
  {$\vect{x}_{k}$};

% Gradient step to the (infeasible) grey point.
\draw[enggray, thick, -{Latex[length=2mm]}] (3.4, 3.65) -- (5.1, 4.13);
\node[enggray, anchor=south, font=\scriptsize] at (4.3, 3.85)
  {gradient step};

% Projection back onto the set: nearest point on the clipped edge.
% The clipped edge runs (4.6,2.6)--(3.2,3.8); project (5.2,4.2) onto it.
% Foot of perpendicular computed by hand: approx (3.86, 3.43).
\fill[engorange] (3.86, 3.43) circle (1.8pt);
\node[engorange, anchor=west, font=\scriptsize] at (3.95, 3.5)
  {$\vect{x}_{k+1} = \Pi_{\mathcal{C}}(\cdot)$};
\draw[engorange, thick, dashed, -{Latex[length=2mm]}] (5.1, 4.13) -- (3.9, 3.46);
\node[engorange, anchor=west, font=\scriptsize] at (4.4, 4.0)
  {projection};

% Right-angle mark at the projection to show orthogonality to the edge.
\draw[engorange, thin]
  (3.98, 3.30) -- (4.08, 3.42) -- (3.96, 3.52);

\node[anchor=north, font=\scriptsize, align=center] at (2.8, -0.75)
  {step, then project onto the nearest feasible point};

\end{tikzpicture}
\caption[The projected-gradient step]{Projected gradient descent turns
a constrained problem into a step-then-project loop. From a feasible
iterate $\vect{x}_{k}$ the method takes an ordinary gradient step,
which may leave the feasible set $\mathcal{C}$ (grey point), then
projects that point back to the nearest feasible point, $\vect{x}_{k+1}
= \Pi_{\mathcal{C}}(\vect{x}_{k} - \eta \grad f(\vect{x}_{k}))$. The
projection onto a convex set is itself a small convex program, and for
the sets an engineer meets most often (a box, a ball, the probability
simplex, a half-space) it has a closed form, so the projection costs
little. The dashed step is orthogonal to the active edge, the
signature of a Euclidean projection, and the method inherits the
convergence rate of unprojected gradient descent whenever the
projection is exact. Projected gradient is the first tool for a convex
problem whose only difficulty is a simple constraint set.}
\label{fig:vol02:ch15:projected-gradient}
\end{figure}
